\section{Contextualização}

O crescimento da utilização da Internet nos últimos anos tem levado os Provedores de Serviços de Internet (ISPs, \textit{Internet Service Providers}) a buscar soluções capazes de garantir a qualidade dos serviços oferecidos aos usuários. O acesso à rede de um ISP, geralmente, constitui uma etapa preliminar para o acesso ao núcleo da Internet \cite{kalmanek2010guide}. Nesse contexto, diferentes usuários e aplicações apresentam requisitos distintos de comunicação, como largura de banda, latência, disponibilidade e confiabilidade.  Dessa forma, os ISPs precisam disponibilizar recursos capazes de atender aos diferentes requisitos estabelecidos para a prestação de seus serviços.

Esses requisitos podem ser estabelecidos por meio de um acordo de nível de serviço (SLA, \textit{Service Level Agreement}), que define as condições acordadas entre o cliente e o ISP para a prestação de determinado serviço. O cumprimento dessas condições está diretamente relacionado à Qualidade de Serviço (QoS, \textit{Quality of Service}) oferecida pelo provedor e pode influenciar a Qualidade de Experiência (QoE, \textit{Quality of Experience}) percebida pelo usuário. Quando os requisitos estabelecidos não são atendidos, a qualidade do serviço pode ser comprometida, resultando em problemas como lentidão e atrasos na comunicação \cite{gomes2020strategies,foukas2017network}.

A necessidade de recursos por parte dos usuários não permanece constante ao longo do tempo. Diferentes períodos podem apresentar variações na quantidade de solicitações e na utilização da infraestrutura, fazendo com que a demanda por recursos da rede aumente ou diminua de acordo com as necessidades dos usuários e das aplicações. Em uma infraestrutura compartilhada, essa característica torna o gerenciamento dos recursos mais complexo, pois uma capacidade dimensionada para atender aos períodos de maior demanda pode permanecer subutilizada em momentos de menor utilização, enquanto uma capacidade insuficiente pode comprometer o atendimento dos usuários durante períodos de maior demanda. Dessa forma, torna-se necessário utilizar os recursos disponíveis de maneira eficiente, adaptando sua distribuição às diferentes condições de demanda.

Nesse contexto, o fatiamento de redes surge como uma abordagem capaz de utilizar uma mesma infraestrutura física para criar diferentes redes virtuais, denominadas fatias, que podem ser configuradas de acordo com os requisitos específicos de cada serviço ou aplicação \cite{vassilaras2017algorithmic}. Dessa forma, os recursos físicos da infraestrutura podem ser compartilhados entre diferentes fatias, permitindo que cada uma possua características e recursos adequados às necessidades dos usuários. Essa flexibilidade possibilita uma utilização mais eficiente da infraestrutura, especialmente em cenários nos quais a demanda por recursos apresenta variações ao longo do tempo.
